% Recuperación autocontenida: sector angular de árbol y reconstrucción Higgs--vacío.
% Propietarios y alcance: higgs_procedencia.md. Incorporar después del operador constitutivo.
\section{Reconstruction conjointe de la jauge, de Higgs et de la réponse constitutive}
\label{amp-higgs-seccion}

La réponse du vide et le secteur électrofaible admettent des lectures complémentaires
du transport angulaire déjà construit. La première conserve les deux canaux
marqués ; le déterminant conserve leur contraction commune ; la composante angulaire
de Higgs conserve en outre l’orientation qui disparaît dans ce déterminant.
La composition de ces lectures permet de reconstruire leur antécédent spectral
et d’exprimer l’impédance au moyen de l’autocouplage scalaire et des couplages
de jauge, dans une même réalisation à l’arbre.

\subsection{Provenance angulaire et signatures des chemins}

On conserve la chaîne APP--TRIT--TPK : le prolongement survivant détermine
le chemin, celui-ci produit son registre orienté et le lecteur du registre exprime
la signature angulaire. Les coordonnées $\pi=\pi_{\rm HMT}$ et
$\alpha=\alpha_{\rm HMT}$ sont des sorties préalables de la même génération ;
$A$ et $C^*$ sont leurs représentations angulaires ultérieures. Dans ce bloc, on
utilise des radians,
\begin{equation}
 x=\frac{\pi A_{\deg}}{180},\qquad
 y=\frac{\pi C^*_{\deg}}{180},\qquad
 A_{\deg}=1000\alpha,\qquad 0<y<x.
 \label{amp-higgs-angulos}
\end{equation}
Par conséquent, $x=50\pi\alpha/9$. Les canaux et les marques de feuillet sont ceux
de la section~\ref{iii:sec:hojas} :
\begin{equation}
 \begin{aligned}
 q_+&=e^{-x-y},& q_-&=e^{-x+y},& D&=q_+q_-=e^{-2x},\\
 P_\pm&=\frac{I\pm\mathcal R}{2},&
 T&=q_+P_++q_-P_-.
 \end{aligned}
 \label{amp-higgs-transporte}
\end{equation}
Ici, $\mathcal R$ est l’involution autoadjointe marquée et chaque projecteur
est de rang un dans la réalisation bicouche de base.

Le lecteur angulaire de durée--perdurance est
\begin{equation}
 \chi_{\rm dur}(n,s)=\exp\!\left(nx+\frac{s}{6}y\right).
 \label{amp-higgs-caracter}
\end{equation}
Pour conserver la provenance finie des coefficients qui seront utilisés,
nous réunissons les empreintes nominales des trois chemins du calendrier de longueur
$108$. Les colonnes de déplacements cardinaux se lisent comme $(N,E,S,O)$ ;
la paire suivante conserve le feuillet et le résidu exprimés par ce chemin.
Le tableau conserve, pour chaque chemin identifié, l’empreinte et la signature
angulaire exprimées par la chaîne précédente.
\begin{center}
\small
\begin{tabular}{c c c c c}
\hline
Secteur & Chemin et axe & $(N,E,S,O)$ & Feuillet/résidu & $(n,s)$\\
\hline
$W$ & $15/37$, E--O & $(23,44,11,30)$ & $(41,8)$ & $(94,-1)$\\
$Z$ & $20/23$, E--O & $(33,45,10,20)$ & $(57,7)$ & $(95,-1)$\\
$H$ & $24/29$, N--S & $(40,34,22,12)$ & $(47,10)$ & $(97,9)$\\
\hline
\end{tabular}
\end{center}
Les trois empreintes typées conservent l’identifiant de chemin, l’orientation
et le registre de provenance des signatures employées. L’évaluation angulaire
commence après l’extraction à partir du registre enrichi ; le tableau
réunit ses sorties pour composer les caractères dans cette réalisation.
Les marques $W_\pm$ qui suivent désignent des coordonnées d’enroulement,
distinctes du nom de l’espèce $W$ :
\begin{equation}
 \begin{aligned}
 W_\pm&=6n\pm s, &
 n&=\frac{W_++W_-}{12},& s&=\frac{W_+-W_-}{2},\\
 (n,s)_W&=(94,-1), & (W_+,W_-)_W&=(563,565),\\
 (n,s)_Z&=(95,-1), & (W_+,W_-)_Z&=(569,571),\\
 (n,s)_H&=(97,9), & (W_+,W_-)_H&=(591,573).
 \end{aligned}
 \label{amp-higgs-firmas}
\end{equation}
La matrice de la première application est
$\bigl(\begin{smallmatrix}6&1\\6&-1\end{smallmatrix}\bigr)$, de déterminant
$-12$ et d’image $\{(u,v)\in\mathbb Z^2:u+v\in12\mathbb Z\}$.
En effet, les formules inverses précédentes sont entières précisément sur
cette image. Le caractère conserve les deux orientations au moyen de
\begin{equation}
 \chi_{\rm dur}(n,s)
 =\left(q_+^{-W_+}q_-^{-W_-}\right)^{1/12}.
 \label{amp-higgs-caracter-canales}
\end{equation}
La preuve s’obtient en prenant les logarithmes : le numérateur est
$W_+(x+y)+W_-(x-y)=12nx+2sy$.
Ainsi sont identifiés la racine douzième et le diviseur $6$ du lecteur.

\begin{proposition}[Quotients angulaires déterminés par les signatures]
Dans une section d’échelle commune, avec les facteurs de raffinement partagés
qui permettent leur annulation dans les quotients, on a
\begin{equation}
 \frac{m_W^{\angle}}{m_Z^{\angle}}=e^{-x},\qquad
 \frac{m_H^{\angle}}{m_W^{\angle}}
 =e^{3x+5y/3},\qquad
 \left(\frac{m_H^{\angle}}{m_W^{\angle}}\right)^2
 =e^{6x+10y/3}.
 \label{amp-higgs-cocientes}
\end{equation}
\end{proposition}
\begin{proof}
Diviser les caractères de $W$ et $Z$ soustrait leurs signatures : $(94,-1)-(95,-1)=(-1,0)$.
Pour $H/W$, la différence est $(97,9)-(94,-1)=(3,10)$.
L’équation~\eqref{amp-higgs-caracter} donne $3x+(10/6)y$ ; le carré
multiplie cet exposant par deux. L’échelle commune et les raffinements
partagés s’annulent avant ces opérations.
\end{proof}
Ces quotients conservent l’identifiant de la coordinatisation angulaire. L’application
à des masses ayant des raffinements sectoriels distincts conserve leurs facteurs
supplémentaires ; la condition d’annulation fait partie de l’énoncé.

\subsection{Réalisation rationalisée à l’arbre}

La lecture électrofaible de cette coordinatisation définit
$c_W^2=D$, $s_W^2=1-D$ et $e_g^2=4\pi\alpha$. La racine positive fixe
$g=e_g/s_W$ et $g'=e_g/c_W$. On obtient
\begin{equation}
 g^2=\frac{4\pi\alpha}{1-D},\qquad
 g'^2=\frac{4\pi\alpha}{D}.
 \label{amp-higgs-calibre}
\end{equation}
La quantité $e_g$ est un couplage adimensionnel dans la normalisation
rationalisée, distinct d’une section dimensionnelle de charge.
La convention du potentiel scalaire est
\begin{equation}
 \begin{gathered}
 V(H)=-\mu_H^2H^\dagger H+\lambda_H(H^\dagger H)^2,\\
 m_H^2=2\lambda_Hv^2,\qquad m_W^2=\frac{g^2v^2}{4}.
 \end{gathered}
 \label{amp-higgs-potencial}
\end{equation}
La même $v$ intervient dans les deux relations. En les divisant, son échelle
s’annule et $m_H^2/m_W^2=8\lambda_H/g^2$. Avec
\eqref{amp-higgs-cocientes}, on obtient
\begin{equation}
 \boxed{\lambda_H=\frac{g^2}{8}\exp\!\left(6x+\frac{10}{3}y\right).}
 \label{amp-higgs-lambda}
\end{equation}
Le facteur $8$ provient des deux normalisations quadratiques de
\eqref{amp-higgs-potencial} ; les coefficients $6$ et $10/3$ proviennent des
signatures. La distinction conserve l’opération qui produit chaque facteur.
La composition reçoit les sorties HMT préalables et utilise cette réalisation physique
à l’arbre à échelle commune ; un changement de normalisation du champ, d’échelle ou
de schéma se transporte également dans ses relations de masses et de couplages.

\subsection{Inverse conjoint et domaine exact}

\begin{theorem}[Reconstruction angulaire au moyen de la jauge et de Higgs]
\label{amp-higgs-inversa-teorema}
Pour $\pi>0$ fixé, l’application de
\eqref{amp-higgs-transporte}, \eqref{amp-higgs-calibre} et
\eqref{amp-higgs-lambda}, sur $\alpha>0$, $x>y>0$, est injective et a
pour inverse
\begin{equation}
 \begin{aligned}
 D&=\frac{g^2}{g^2+g'^2},&
 x&=-\frac12\log D,&
 \alpha&=\frac{g^2g'^2}{4\pi(g^2+g'^2)},\\
 y&=\frac3{10}\left[\log\!\left(\frac{8\lambda_H}{g^2}\right)-6x\right].
 \end{aligned}
 \label{amp-higgs-inversa}
\end{equation}
Son image exacte est formée des triplets $g,g',\lambda_H>0$ tels que,
avec $D$ et $x$ reconstruits comme ci-dessus,
\begin{equation}
 \frac{g^2}{8}e^{6x}<\lambda_H<\frac{g^2}{8}e^{28x/3}.
 \label{amp-higgs-dominio}
\end{equation}
\end{theorem}
\begin{proof}
Les deux équations de jauge donnent
$g^2/(g^2+g'^2)=D$ et
$g^2g'^2/(g^2+g'^2)=4\pi\alpha$. Le logarithme réel retrouve $x$ et
l’équation scalaire retrouve $y$. Les conditions $g,g'>0$ impliquent
$0<D<1$ et $x>0$. L’inégalité~\eqref{amp-higgs-dominio} équivaut,
par le logarithme strictement croissant, à
$0<\log(8\lambda_H/g^2)-6x<10x/3$, c’est-à-dire $0<y<x$.
Réciproquement, les formules reconstruites produisent
$4\pi\alpha/(1-D)=g^2$, $4\pi\alpha/D=g'^2$ et
$g^2e^{6x+10y/3}/8=\lambda_H$. Les racines positives retrouvent les
couplages originaux. Cela prouve la nécessité, la suffisance et l’unicité.
\end{proof}

L’inversion précédente décrit exactement cette application entre lecteurs.
La section engendrée par HMT conserve en outre la relation de
\eqref{amp-higgs-angulos}, qui s’exprime comme
\begin{equation}
 -\frac12\log\!\left(\frac{g^2}{g^2+g'^2}\right)
 =\frac{25}{18}\frac{g^2g'^2}{g^2+g'^2}.
 \label{amp-higgs-compatibilidad-alpha}
\end{equation}
Elle s’obtient en substituant l’inverse de $\alpha$ dans $x=50\pi\alpha/9$.
Le contre-angle conserve également son équation de retour antérieure.
Par conséquent, l’image du lecteur et la section produite par le générateur
sont des objets distincts : les compatibilités préalables sélectionnent la section
au sein du domaine reconstructible. L’opération inverse agit après
la génération et permet de vérifier ses sorties conservées.

\subsection{Reconstruction Higgs--impédance}

Pour exprimer le contenu orienté sans altérer les canaux,
nous définissons
\begin{equation}
 \rho_{\rm or}=\frac{q_-}{q_+}=e^{2y},\qquad
 \Xi_H=\frac{8D^3\lambda_H}{g^2}.
 \label{amp-higgs-Xi}
\end{equation}
L’équation~\eqref{amp-higgs-lambda} et $D^3=e^{-6x}$ donnent
\begin{equation}
 \begin{aligned}
 \Xi_H&=e^{10y/3}=\rho_{\rm or}^{5/3},&
 \rho_{\rm or}&=\Xi_H^{3/5},\\
 q_-&=\sqrt D\,\Xi_H^{3/10},&
 q_+&=\sqrt D\,\Xi_H^{-3/10}.
 \end{aligned}
 \label{amp-higgs-canales}
\end{equation}
Les puissances sont les racines réelles positives. Le domaine exact
de la chambre orientée est
\begin{equation}
 0<D<1,\qquad 1<\Xi_H<D^{-5/3}.
 \label{amp-higgs-dominio-Xi}
\end{equation}
En effet, $0<y<x$ équivaut à $1<e^{10y/3}<e^{10x/3}$.
L’extrémité inférieure donne $q_+=q_-$ ; la supérieure donne $q_-=1$.
Pour les deux orientations ensemble, $|y|<x$ équivaut à
$D^{5/3}<\Xi_H<D^{-5/3}$.

Soit $F(q)=f(q^{30})=(1-q^{90})/(1-q^{120})$, la réponse déjà prouvée dans
\eqref{iii:eq:rchannels}--\eqref{iii:eq:monotone}.
\begin{corollary}[Confluence des lectures scalaire et constitutive]
\label{amp-higgs-impedancia-corolario}
Dans le domaine~\eqref{amp-higgs-dominio-Xi},
\begin{equation}
 \boxed{\widehat Z=
 \frac{F(\sqrt D\,\Xi_H^{3/10})}
 {F(\sqrt D\,\Xi_H^{-3/10})}.}
 \label{amp-higgs-impedancia}
\end{equation}
Les réponses complètes sont
\begin{equation}
 \begin{aligned}
 r_-&=F(\sqrt D\,\Xi_H^{3/10}),&
 r_+&=F(\sqrt D\,\Xi_H^{-3/10}),\\
 \widehat\mu&=r_-^2,& \widehat\varepsilon&=r_+^2,&
 \widehat c&=(r_-r_+)^{-1}.
 \end{aligned}
 \label{amp-higgs-vacio}
\end{equation}
De plus, les deux reconstructions satisfont
\begin{equation}
 \frac{F^{-1}(\sqrt{\widehat\mu})}
 {F^{-1}(\sqrt{\widehat\varepsilon})}
 =\rho_{\rm or}=\Xi_H^{3/5}.
 \label{amp-higgs-confluencia}
\end{equation}
\end{corollary}
\begin{proof}
Substituer les canaux de~\eqref{amp-higgs-canales} dans les quatre lectures
de~\eqref{iii:eq:four} prouve les premières identités. L’inverse
constitutif de~\eqref{iii:eq:cubic} retrouve $q_-$ et $q_+$ ; leur quotient
est $\rho_{\rm or}$. L’égalité avec $\Xi_H^{3/5}$ est déjà prouvée dans
\eqref{amp-higgs-canales}.
\end{proof}

Dans la section compatible avec $A_{\deg}=1000\alpha$, le couplage
$Q(D):=-9\log D/25$ satisfait $Q(D)=4\pi\alpha$ et
$g^2=Q(D)/(1-D)$. Par conséquent, la même coordonnée de Higgs s’écrit
\begin{equation}
 \Xi_H=\frac{8(1-D)D^3\lambda_H}{Q(D)},\qquad
 \frac1{g^2}+\frac1{g'^2}=\frac1{Q(D)}.
 \label{amp-higgs-curva}
\end{equation}
Ces expressions relient le couplage total, sa répartition rationalisée
et l’orientation scalaire dans la même section.

\subsection{Information conservée et transport d’orientation}

La paire de jauge retrouve $\alpha$ et $x$ ; le quotient
$\lambda_H/g^2$ retrouve ensuite $y$. Avec les marques $P_\pm$, on reconstruit
les opérateurs complets $T$ et $\mathcal V=r_+P_++r_-P_-$.
Les trois scalaires déterminent les valeurs propres ; les projecteurs marqués
déterminent en outre leur réalisation opératorielle. L’histoire enrichie du
chemin reste une information antérieure de résolution supérieure.

Sous $y\mapsto-y$ par rapport aux marques fixes, $D,g,g'$ restent
invariants et $\Xi_H\mapsto\Xi_H^{-1}$. Les réponses constitutives s’échangent,
$\widehat Z\mapsto\widehat Z^{-1}$ et $\widehat c$ est conservé.
Si $\lambda_H^+$ et $\lambda_H^-$ désignent les évaluations dans les deux
orientations conjuguées, on a
\begin{equation}
 \lambda_H^+\lambda_H^-
 =\frac{g^4}{64}D^{-6}.
 \label{amp-higgs-orientacion}
\end{equation}
La preuve consiste à multiplier~\eqref{amp-higgs-lambda} pour $y$ et $-y$ ;
les contributions orientées s’annulent. Cette comparaison conserve le
lecteur de signatures et change son argument angulaire. Une conjugaison simultanée
des chemins et des références doit également transporter les signatures.

Dans l’amplification $T_m=T\otimes I_{9^m}$, avec
$\tau_m=\operatorname{Tr}/(2\cdot9^m)$ et
$\mathcal R_m=\mathcal R\otimes I_{9^m}$, les coordonnées appropriées sont
\begin{equation}
 D=\exp\bigl(2\tau_m(\log T_m)\bigr),\qquad
 \rho_{\rm or}=\exp\bigl(-2\tau_m(\mathcal R_m\log T_m)\bigr).
 \label{amp-higgs-refinamiento}
\end{equation}
En effet, $\log T_m=(\log T)\otimes I_{9^m}$, de sorte que les deux traces
sont $-x$ et $-y$, respectivement. La normalisation élimine la multiplicité
de l’amplification. Les lectures de jauge, de Higgs et du vide restent ainsi
compatibles avec ces raffinements ; le déterminant ordinaire amplifié
serait $D^{9^m}$ et a une fonction distincte.

La conclusion de la composition est précise : l’invariant pair de jauge
et la coordonnée orientée scalaire reconstruisent conjointement la paire
spectrale qui gouverne le vide. Lorsque $\lambda_H$ est évalué au moyen de la
même formule, la confluence est une identité de cohérence. Une lecture
indépendante de $\lambda_H$, transportée vers des schéma et échelle identiques,
transforme cette égalité en comparaison supplémentaire. Dans les deux usages restent
explicites les signatures, les marques de feuillet et la normalisation à l’arbre.
